\documentclass[a4paper,10pt, oneside]{book}
%\usepackage[latin1]{inputenc}
\usepackage[spanish]{babel} % Normas tipográficas y opciones del español e ingles
\usepackage[T1]{fontenc}    		% Codificación de salida
\usepackage[utf8]{inputenc} 		% Codificación de entrada 
%\usepackage{subfig}
%\def\spanishoptions{argentina}
%\usepackage[spanish]{babel}
\usepackage{calc}
\usepackage{fancyhdr}
\usepackage{graphicx}
\usepackage{graphics}

\usepackage{amsmath}
\usepackage{amsthm}
\usepackage{amsfonts}
\usepackage{amssymb}
\usepackage{mathcomp}

\usepackage{fancyhdr,fancybox} %%Fancy headings
\usepackage{multirow}
\usepackage{hyperref}
\usepackage{subfigure}
\usepackage{ccaption}
\usepackage{enumerate}

\begin{document}

\section{Metodología}

Dada la limitación del personal disponible encargada de las tareas del proyecto, se planifica su realización de forma secuencial con relación terminar para empezar (TE) y la utilización de una metodología del tipo cascada ya que los requerimientos son inmutables. Separando el proyecto en tres etapas:

\begin{itemize}
\item Diseño y desarrollo de los métodos.
\item Implementación y Análisis de los métodos.
\item Cierre
\end{itemize}

Cada fase finalizará con la documentación de las tareas realizadas en base a criterios objetivos que así lo determinen. %TODO agregar algunas cuestiones referidas a estos aca


\section{Plan de Tareas}


\section{Cronograma}


\end{document}